% Interior design 5e: the two box families. Every environment of the manuscripts belongs
% to one of them (the full mapping is in docs/design/interior.md).
%   Formal   theorem, proposition, lemma, corollary, definition, example and their unnumbered
%            x-variants: an 8% charcoal fill under one 1.2 pt coral rule, never split; the
%            statement is upright, the label is a Fira capital label with a coral number and
%            an EB Garamond italic note.
%   Teaching exercise, reminder, think (Pause and think), beauty (Why this is beautiful),
%            challenge, remark: a 2 pt charcoal bar on the left, no fill; may run on to the
%            next page. Proofs, hints and solutions are the run-in members of the family:
%            the same labels, no bar, so long arguments flow with the text.
\makeatletter
% ---------------------------------------------------------------- formal family
\colorlet{ocFill}{ocBlack!8}
\newtheoremstyle{ocformal}{0pt}{0pt}{\normalfont}{}{}{}{0.8em}
 {{\ocLabelFont\ocLab{\thmname{#1}}\thmnumber{\ \ocNum{#2}}}\thmnote{\enspace\ocNote{#3}}}
\theoremstyle{ocformal}
\newtheorem{theorem}{Theorem}[chapter]
\newtheorem{proposition}[theorem]{Proposition}
\newtheorem{lemma}[theorem]{Lemma}
\newtheorem{corollary}[theorem]{Corollary}
\newtheorem{definition}[theorem]{Definition}
\newtheorem{example}[theorem]{Example}
% Unnumbered variants, so a statement or worked calculation does not shift the numbering.
\newtheorem*{theoremx}{Theorem}
\newtheorem*{propositionx}{Proposition}
\newtheorem*{lemmax}{Lemma}
\newtheorem*{corollaryx}{Corollary}
\newtheorem*{definitionx}{Definition}
\newtheorem*{examplex}{Example}
\tcbset{oc formal/.style={enhanced,unbreakable,arc=0pt,frame hidden,boxrule=0pt,colback=ocFill,
   borderline north={1.2pt}{0pt}{ocCoral},left=10pt,right=10pt,top=8pt,bottom=8pt,boxsep=0pt,
   before skip=13pt plus 3pt,after skip=13pt plus 3pt}}
\@for\oc@env:={theorem,proposition,lemma,corollary,definition,example,theoremx,propositionx,%
  lemmax,corollaryx,definitionx,examplex}\do{\expandafter\tcolorboxenvironment\expandafter{\oc@env}{oc formal}}
% ---------------------------------------------------------------- teaching family
\tcbset{oc teach/.style={enhanced,breakable,arc=0pt,frame hidden,boxrule=0pt,interior hidden,
   borderline west={2pt}{0pt}{ocBlack},left=12pt,right=0pt,top=1pt,bottom=1pt,boxsep=0pt,
   before skip=14pt plus 3pt,after skip=14pt plus 3pt,lines before break=3,
   % Teaching boxes may continue on the next page: the bar simply continues.
   no finish first,no finish middle,no finish last},
  oc title/.style={before upper={{\ocLabelFont #1}\par\vspace{4pt}}}}
\newcommand{\ocBoxTitle}[2]{\ocLab{#1}\ifx\relax#2\relax\else\enspace\ocNote{#2}\fi}
\newtcolorbox{reminder}[1][]{oc teach,oc title={\ocBoxTitle{Reminder}{#1}}}
\newtcolorbox{beauty}{oc teach,oc title={\ocBoxTitle{Why this is beautiful}{}}}
\newtcolorbox{think}{oc teach,oc title={\ocBoxTitle{Pause and think}{}}}
\newtcolorbox{challenge}{oc teach,oc title={\ocBoxTitle{Challenges}{}}}
\newtcolorbox{remark}[1][]{oc teach,before upper={{\ocLabelFont\ocLab{Remark}}%
  \ifx\relax#1\relax\else\enspace\ocNote{#1}\fi\hspace{0.8em}\ignorespaces}}
% Exercises and solutions. The counter anchor travels inside the box.
\newcounter{exo}[chapter]
% The exercise name. A problem chapter switches it with \renewcommand{\ocExName}{Problem}.
\newcommand{\ocExName}{Exercise}
\renewcommand{\theexo}{\thechapter.\arabic{exo}}
\newtcolorbox{exercise}[1][]{oc teach,top=2pt,bottom=2pt,before skip=12pt plus 3pt,after skip=12pt plus 3pt,
  before upper={\refstepcounter{exo}\parindent0pt{\ocLabelFont\ocLab{\ocExName}\ \ocNum{\theexo}}%
    \ifx\relax#1\relax\else\enspace{\color{ocBlack}#1}\fi\hspace{0.8em}\ignorespaces}}
% Hints and solutions run in, with the same charcoal label and coral number.
\newenvironment{hint}[1]{%
  \par\ocOpeningSpace\smallskip\noindent{\ocLabelFont\ocLab{Hint for \ocExName}\ \ocNum{\ref{#1}}}\hspace{0.8em}\ignorespaces}%
  {\par\smallskip}
\newenvironment{solution}[1]{%
  \par\ocOpeningSpace\medskip\noindent{\ocLabelFont\ocLab{Solution of \ocExName}\ \ocNum{\ref{#1}}}\hspace{0.8em}\ignorespaces}%
  {\hfill\ocDia{ocCoral}\par\medskip}
% Solutions: \idea explains how to find the approach, \writeup gives the full argument.
\newcommand{\idea}{\par\nobreak\noindent\textit{How to find it.}\quad}
\newcommand{\writeup}{\par\smallskip\noindent\textit{Write-up.}\quad}
\newcommand{\tools}[1]{\par\smallskip\noindent\textit{Tools used.}\quad #1}
\newcommand{\leads}{\par\smallskip\noindent\textit{Where it leads.}\quad}
% Discarded in the PDF. The EPUB reader uses the same environment as a name switch.
\newbox\ocExNameBox
\newenvironment{exname}{\setbox\ocExNameBox\vbox\bgroup}{\egroup}
% Proofs: italic run-in label, the coral diamond at the end.
\renewcommand{\qedsymbol}{\ocDia{ocCoral}}
\renewenvironment{proof}[1][\proofname]{\par\pushQED{\qed}\normalfont\topsep6pt\relax
  \trivlist\item[\hskip\labelsep{\itshape #1.}]\ignorespaces}{\popQED\endtrivlist\@endpefalse}
\makeatother
